% Guide utilisateur complet d'EP Manager (mod_ep et mod_epsynthesis) : tous les rôles.
% Compilation : latexmk -pdf guide-utilisateur.tex
% Captures : node captures.js <site.json> gu-... (voir l'annexe « Captures »).
\documentclass[11pt,a4paper]{article}

% Version du guide : à mettre à jour à chaque modification, avec une ligne dans l'historique.
\newcommand{\titreguide}{Guide utilisateur}
\newcommand{\versionguide}{1.2}
\newcommand{\dateguide}{novembre 2026}
\newcommand{\versionep}{1.0.5}
\newcommand{\versionsynthese}{1.0.2}

\usepackage{guide-commun}
\setcounter{tocdepth}{2}

% Étape numérotée d'une procédure, en marge de la liste.
\newlist{etapes}{enumerate}{1}
\setlist[etapes]{label=\textbf{\arabic*.}, leftmargin=*, itemsep=3pt}

\begin{document}

\pagedegarde{Étudiants, enseignants, DEVE : l'enseignement personnalisé de bout en bout, pas à
pas}{gu-etu-tableau-de-bord}

\tableofcontents

\section*{Historique des versions}
\addcontentsline{toc}{section}{Historique des versions}

\begin{center}
\begin{tabularx}{\linewidth}{@{}l l X@{}}
\toprule
\textbf{Version} & \textbf{Date} & \textbf{Modifications} \\
\midrule
1.0 & novembre 2026 & Première version, pour EP Manager 1.0.3 : tous les rôles en un seul
document, captures prises sur une promotion de 1\,000 étudiants. \\
1.1 & novembre 2026 & Séances du programme sans date (EP Manager 1.0.4) ; qui valide les
syllabus. \\
1.2 & novembre 2026 & EP de stage recalculés dès qu'un stage change (EP Manager 1.0.5). \\
\bottomrule
\end{tabularx}
\end{center}

\section*{Comment lire ce guide}
\addcontentsline{toc}{section}{Comment lire ce guide}

Ce guide décrit tout EP Manager, pour tous ceux qui s'en servent. Chacun peut s'en tenir à sa
partie :
\begin{center}
\begin{tabularx}{\linewidth}{@{}l X r@{}}
\toprule
\textbf{Vous êtes} & \textbf{Lisez} & \textbf{Page} \\
\midrule
Tout le monde & Partie~\ref{part:principes} : principes, vocabulaire, statuts, calendrier &
\pageref{part:principes} \\
Étudiant & Partie~\ref{part:etudiant} & \pageref{part:etudiant} \\
Enseignant (responsable d'EP, référent, auteur) & Partie~\ref{part:enseignant} &
\pageref{part:enseignant} \\
DEVE & Partie~\ref{part:deve} & \pageref{part:deve} \\
Administrateur du site Moodle & Partie~\ref{part:admin} & \pageref{part:admin} \\
Tout le monde, en cas de doute & Partie~\ref{part:faq} : questions fréquentes &
\pageref{part:faq} \\
\bottomrule
\end{tabularx}
\end{center}

Les captures d'écran ont été prises sur le banc d'essai, dans une promotion de 1\,000 étudiants,
150 enseignants et 60 EP au catalogue, après que tous les étudiants se sont inscrits, ont
déclaré des EP et demandé des éléments de leur supplément, que leurs enseignants et la DEVE ont
tout tranché, et que les suppléments ont été émis (voir le rapport \emph{Test des interactions
à 1\,000 étudiants}). Les noms sont fictifs.

Conventions : \chemin{Administration > Types d'EP} désigne un chemin de navigation dans
l'interface, \bouton{Enregistrer} un bouton ou un lien, \statut{En attente} un statut affiché.

\newpage

% ===========================================================================
\part{Principes}\label{part:principes}
% ===========================================================================

\section{Ce que fait EP Manager}

Les enseignements personnalisés (EP) sont les crédits ECTS que chaque étudiant doit réunir, au
fil de son cursus, en plus du tronc commun : EP du catalogue de l'école, stages
complémentaires, engagement étudiant, expérience professionnelle, sport, enseignement suivi
dans un autre établissement\dots{} EP Manager tient le dossier de chaque étudiant, de
l'inscription ou de la déclaration jusqu'au supplément au diplôme :
\begin{itemize}
  \item l'étudiant s'inscrit aux EP du catalogue et déclare ce qu'il a fait ailleurs, avec ses
    justificatifs ;
  \item chaque demande suit le \textbf{circuit de validation} de son type : responsable de
    l'EP, enseignant référent, DEVE, ou plusieurs d'entre eux l'un après l'autre ;
  \item les ECTS validés sont comptés face aux \textbf{minimums} annuels et du cursus, dans la
    limite des \textbf{plafonds} de chaque type ;
  \item en fin d'année, la DEVE valide les années d'étude ; en fin de cursus, elle émet le
    \textbf{supplément au diplôme} certifié de chaque étudiant.
\end{itemize}

\subsection{Deux activités Moodle}

\begin{description}
  \item[Enseignement personnalisé] une activité par promotion, dans le cours de la promotion
    (celui qui porte déjà \og Gestion des stages\fg{}). C'est là que tout se passe : dossiers
    des étudiants, catalogue, décisions, paramétrage.
  \item[Suivi de l'enseignement personnalisé] une activité dans un cours d'équipe, qui
    rassemble pour chaque enseignant ce qu'il a à faire dans \emph{toutes} les promotions. Elle
    n'accorde aucun droit : elle montre et renvoie vers les activités des promotions.
\end{description}

\subsection{Qui fait quoi}

\begin{center}
\begin{tabularx}{\linewidth}{@{}l X@{}}
\toprule
\textbf{Rôle} & \textbf{Ce qu'il fait} \\
\midrule
Étudiant & S'inscrit aux EP du catalogue, déclare les autres, joint ses justificatifs, suit ses
décisions et ses ECTS, demande des éléments pour son supplément, télécharge son supplément
émis. \\
Responsable d'EP & Enseignant désigné pour un EP du catalogue : inscrit ou retire des
étudiants, note ou valide les inscrits (un par un, en lot ou par import CSV), clôt l'EP. \\
Enseignant référent & Celui qui suit l'étudiant dans \og Gestion des stages\fg{} : tranche les
déclarations de ses étudiants et leurs éléments du supplément, à son étape du circuit. \\
Auteur & Tout enseignant : propose un nouvel EP, rédige son syllabus, met à jour celui des EP
dont il est responsable. \\
DEVE & Paramètre tout, tranche et corrige toute décision, examine les propositions, pilote la
promotion, valide les années, émet les suppléments, importe et exporte. \\
\bottomrule
\end{tabularx}
\end{center}

Chacun ne voit que ce qui le concerne : un étudiant, son dossier ; un enseignant, ses EP et ses
étudiants référents ; la DEVE, toute la promotion.

\section{Vocabulaire}\label{sec:vocabulaire}

\begin{description}
  \item[EP] enseignement personnalisé : une ligne du dossier de l'étudiant, avec un type, une
    année d'étude, des ECTS et un statut.
  \item[Type d'EP] catégorie qui fixe les règles : comment l'EP arrive au dossier (mode
    d'acquisition), qui le valide (circuit), s'il est noté, quels justificatifs il faut, quels
    plafonds s'appliquent. Six types sont prévus d'office : académique interne, stage,
    engagement étudiant, expérience professionnelle, sport, académique externe.
  \item[Catalogue] les EP proposés par l'école, auxquels l'étudiant s'inscrit. Chacun a sa
    fiche (\textbf{syllabus}), ses places, sa période d'inscription et ses responsables.
  \item[Déclaration] EP suivi hors du catalogue, déclaré par l'étudiant avec ses justificatifs.
  \item[Circuit de validation] suite d'une à trois étapes (responsable de l'EP, enseignant
    référent, DEVE) par lesquelles passe une demande ; une demande validée à une étape passe
    à la suivante ; refusée à n'importe quelle étape, elle s'arrête.
  \item[ECTS demandés, retenus, validés] l'étudiant \emph{demande} des ECTS ; celui qui valide
    \emph{retient} un nombre d'ECTS (au plus celui demandé) ; les ECTS \emph{validés} d'un type
    ne \emph{comptent} que dans la limite de ses plafonds.
  \item[Plafond] maximum d'ECTS d'un type pris en compte, sur le cursus ou par année. Au-delà,
    les ECTS restent validés mais sont \og non retenus\fg{}.
  \item[Minimum] nombre d'ECTS à réunir pour une année d'étude, ou sur l'ensemble du cursus.
  \item[Élément complémentaire] activité que l'étudiant demande à faire figurer dans son
    supplément au diplôme (rubrique 6.1) : responsabilité associative, prix, mobilité,
    certification\dots
  \item[Supplément au diplôme] document au modèle européen Europass qui décrit le diplôme et
    le parcours de l'étudiant ; émis par la DEVE, numéroté et vérifiable.
\end{description}

\section{Statuts d'une demande}\label{sec:statuts}

\begin{center}
\begin{tabularx}{\linewidth}{@{}l X@{}}
\toprule
\textbf{Statut} & \textbf{Signification} \\
\midrule
\statut{En attente} & La demande attend une décision ; la page indique qui doit décider
(\og en attente : Enseignant référent\fg{}, par exemple). L'étudiant peut encore la modifier
(déclaration) ou la retirer. \\
\statut{Validé} & Décision finale favorable : les ECTS retenus (et la note, pour un EP noté)
comptent dans le bilan. \\
\statut{Refusé} & Décision défavorable, à n'importe quelle étape, avec le commentaire de celui
qui a décidé. Pour un EP noté, une note sous le seuil vaut refus. \\
\statut{Retiré} & Retirée par l'étudiant (tant qu'elle était en attente), par le responsable
ou par la DEVE. Une inscription retirée rend sa place. \\
\bottomrule
\end{tabularx}
\end{center}

Chaque décision est gardée dans l'\textbf{historique} de la demande : qui, quand, à quel titre
(responsable, référent, DEVE), avec quelle note, quels ECTS et quel commentaire. Deux personnes
qui tranchent la même demande au même moment ne peuvent pas l'enregistrer deux fois : la
seconde est prévenue que la demande a changé, et rien n'est écrit.

\section{Messages envoyés}\label{sec:messages}

Les messages partent par les canaux choisis par chacun dans ses préférences Moodle
(notification, courriel). Leur texte est réglé par la DEVE (section~\ref{sec:courriels}).

\begin{center}
\begin{tabularx}{\linewidth}{@{}X l@{}}
\toprule
\textbf{Message} & \textbf{Destinataire} \\
\midrule
Décision finale sur une demande d'EP & L'étudiant \\
Proposition d'EP (ou mise à jour de syllabus) à examiner & La DEVE \\
Décision sur une proposition & L'enseignant auteur \\
Relance hebdomadaire (le lundi) : demandes en attente depuis plus d'une semaine & Chaque
validateur concerné \\
Fin de l'émission groupée des suppléments & La DEVE qui l'a lancée \\
\bottomrule
\end{tabularx}
\end{center}

\section{Une année avec EP Manager}\label{sec:calendrier}

\begin{center}
\begin{tabularx}{\linewidth}{@{}l X X@{}}
\toprule
\textbf{Moment} & \textbf{DEVE} & \textbf{Enseignants et étudiants} \\
\midrule
Été & Rentrée de l'activité (année suivante), règles copiées ou ajustées, catalogue mis à
jour & Les auteurs proposent leurs nouveaux EP et mettent à jour leurs syllabus \\
Rentrée & Publication du catalogue, périodes d'inscription & Les étudiants s'inscrivent \\
Année & Pilotage, relances, décisions de dernière étape & Les étudiants déclarent ; les
responsables notent leurs EP à la fin de chacun ; les référents tranchent les déclarations \\
Fin d'année & Validation des années (jury) & Les étudiants vérifient leur bilan \\
Fin de cursus & Éléments du supplément, émission des suppléments & Les étudiants demandent
leurs éléments complémentaires, téléchargent leur supplément \\
\bottomrule
\end{tabularx}
\end{center}

\newpage
% ===========================================================================
\part{Étudiant}\label{part:etudiant}
% ===========================================================================

\section{Votre tableau de bord}

Ouvrez le cours de votre promotion, puis l'activité \og Enseignement personnalisé\fg{}. Vous
arrivez sur \textbf{Mes enseignements personnalisés}.

\capture{gu-etu-tableau-de-bord}{Tableau de bord d'un étudiant en fin d'année : années, types,
détail des EP.}

Il se lit de haut en bas :
\begin{description}
  \item[Synthèse] le total d'ECTS retenus, le nombre d'années validées et les ECTS encore en
    attente de validation.
  \item[Bilan par année d'étude] pour chaque année, le minimum exigé, les ECTS retenus, ce qui
    reste à valider, et l'état \statut{Atteint} ou \statut{À compléter} ; le minimum de
    cursus s'il y en a un.
  \item[Bilan par type] pour chaque type, ses plafonds, les ECTS validés, retenus, non retenus
    (au-delà d'un plafond) et en attente.
  \item[Détail des enseignements personnalisés] tous vos EP, avec leur statut et leurs
    actions : \bouton{Détail}, \bouton{Modifier} (une déclaration en attente),
    \bouton{Retirer ma demande} (une demande en attente).
\end{description}

La barre de liens en haut de page mène au \chemin{Catalogue des EP}, à \chemin{Déclarer un
EP} et à \chemin{Mon supplément au diplôme}.

\begin{astuce}
Seuls les ECTS \textbf{validés} comptent, et seulement dans la limite des plafonds de leur
type. Des ECTS non retenus restent acquis : ils reviendraient au décompte si le plafond était
relevé.
\end{astuce}

\section{S'inscrire à un EP du catalogue}\label{sec:inscription}

\subsection{Trouver un EP}

Le \chemin{Catalogue des EP} liste les EP ouverts à votre promotion, 40 par page, dans l'ordre
choisi par la DEVE puis par intitulé. La zone de recherche filtre sur l'intitulé, le code et la
description. Chaque ligne indique le type,
les ECTS, les années concernées, les places restantes, les responsables et la période
d'inscription.

\capture{gu-etu-catalogue}{Catalogue des EP : recherche, places, inscription directe.}

Le lien \bouton{Syllabus} ouvre la fiche de l'EP : objectifs, acquis d'apprentissage,
prérequis, contenu, méthodes, conditions d'examen et de seconde session, bibliographie,
programme des séances avec leurs intervenants, volume horaire. Elle se télécharge en PDF.

\capture{gu-etu-syllabus}{Fiche (syllabus) d'un EP du catalogue.}

\capture[0.85\linewidth]{gu-etu-seances}{Programme des séances, tel qu'il figure sur la fiche.}

\subsection{S'inscrire}

\begin{etapes}
  \item Sur la ligne de l'EP, choisissez l'\textbf{année d'étude} au titre de laquelle il
    comptera (votre année courante est proposée).
  \item Cliquez sur \bouton{S'inscrire}.
\end{etapes}

L'inscription est enregistrée aussitôt, \statut{En attente} de la décision du responsable de
l'EP ; vous restez sur la même page du catalogue pour vous inscrire à d'autres EP. Les ECTS ne
comptent qu'une fois l'inscription validée, à la fin de l'EP.

Quand l'inscription n'est pas possible, le bouton est remplacé par la raison :
\begin{itemize}
  \item \statut{Inscriptions fermées} : hors de la période d'inscription, ou EP clos ;
  \item \statut{Complet} : plus de place (deux étudiants qui prennent la dernière place à la
    même seconde ne peuvent pas l'obtenir tous les deux) ;
  \item l'EP n'est pas ouvert à votre année d'étude ;
  \item vous êtes déjà inscrit ;
  \item \statut{Prérequis manquants} : l'EP exige d'avoir validé un autre EP ou un nombre
    minimum d'ECTS. Si ces prérequis sont \textbf{bloquants}, l'inscription est refusée.
\end{itemize}

Des prérequis \textbf{non bloquants} qui manquent sont signalés en orange sous l'intitulé de
l'EP : l'inscription reste possible, et le responsable voit le manque dans sa liste de
participants.

\subsection{Retirer une inscription}

Tant que l'inscription est \statut{En attente}, \bouton{Retirer ma demande} sur votre tableau
de bord la retire et rend la place. Une fois l'inscription tranchée, adressez-vous au
responsable de l'EP ou à la DEVE.

\section{Déclarer un EP hors catalogue}\label{sec:declarer}

Engagement associatif, expérience professionnelle, sport, enseignement suivi dans un autre
établissement\dots{} se déclarent depuis \chemin{Déclarer un EP}.

\capture{gu-etu-declarer}{Déclaration d'un EP : type, intitulé, année, ECTS, justificatifs.}

\begin{etapes}
  \item Lisez le bandeau : il indique à qui la demande sera transmise (votre enseignant
    référent, la DEVE, ou l'un puis l'autre selon le type).
  \item Choisissez le \textbf{type} : sous la liste s'affichent, pour chaque type, son circuit
    de validation, sa consigne et ses plafonds.
  \item Donnez un \textbf{intitulé} et une \textbf{description} : ce que vous avez fait, où,
    combien de temps, avec quel résultat.
  \item Choisissez l'\textbf{année d'étude} et le nombre d'\textbf{ECTS demandés}.
  \item Déposez vos \textbf{justificatifs} (attestation, convention, certificat\dots) par
    glisser-déposer ou avec le sélecteur de fichiers ; ils sont obligatoires pour certains
    types, et les formats acceptés sont rappelés sous la zone de dépôt.
  \item Cliquez sur \bouton{Envoyer la demande}.
\end{etapes}

Votre enseignant référent est celui qui vous est attribué dans l'activité \og Gestion des
stages\fg{} du même cours. S'il n'y en a pas, la page le signale : adressez-vous à la DEVE.

\begin{attention}
Une demande tranchée ne se modifie plus. Tant qu'elle est en attente, \bouton{Modifier} et
\bouton{Retirer ma demande} restent possibles depuis le tableau de bord.
\end{attention}

\subsection{EP de type stage}

Les stages complémentaires validés par la DEVE dans \og Gestion des stages\fg{} donnent
automatiquement des ECTS d'EP, d'après leur durée : vous n'avez rien à déclarer. Ils
apparaissent dans votre tableau de bord dès la validation du stage.

\section{Suivre vos demandes}

Le \bouton{Détail} d'un EP montre la demande telle que vous l'avez faite, vos justificatifs,
l'étape où elle se trouve et l'\textbf{historique des décisions}.

\capture{gu-etu-detail}{Détail d'un EP validé : demande, justificatifs, historique des
décisions.}

Vous recevez un message à chaque décision finale (validé ou refusé), avec le commentaire de
celui qui a décidé. Un refus n'empêche pas de déposer une nouvelle demande mieux étayée.

\section{Votre supplément au diplôme}\label{sec:etu-supplement}

\subsection{Ce que montre la page}

\chemin{Mon supplément au diplôme} réunit :
\begin{itemize}
  \item un bouton \bouton{Aperçu (non certifié)}, en français ou en anglais : le document tel
    qu'il serait émis aujourd'hui, marqué \og APERÇU -- NON CERTIFIÉ\fg{}, avec la liste des
    informations qui manquent encore (date de naissance, signataire\dots) ;
  \item les EP validés qui figureront à la rubrique~4.3, avec leurs ECTS et leur note ;
  \item vos \textbf{éléments ajoutés au supplément} (rubrique~6.1), avec leur statut ;
  \item les \textbf{suppléments émis} par la DEVE, à télécharger.
\end{itemize}

\capture{gu-etu-supplement}{Mon supplément au diplôme : EP retenus, éléments validés,
supplément émis.}

\subsection{Demander l'ajout d'un élément}

\begin{etapes}
  \item Cliquez sur \bouton{Demander l'ajout d'un élément}.
  \item Choisissez la \textbf{catégorie} (responsabilités, mobilité, prix, certification\dots,
    selon la liste de la DEVE).
  \item Indiquez la \textbf{période} (début et fin).
  \item Donnez l'\textbf{intitulé} en français et, si possible, en anglais : c'est lui qui
    figurera dans le supplément édité en anglais.
  \item Décrivez l'activité.
  \item Déposez vos \textbf{preuves} en PDF (facultatives ou obligatoires selon le choix de la
    DEVE).
  \item Enregistrez.
\end{etapes}

\capture{gu-etu-supplement-element}{Demande d'un élément complémentaire du supplément.}

La demande suit le circuit fixé par la DEVE (par exemple votre enseignant référent, puis la
DEVE). Tant qu'elle est en attente, vous pouvez la modifier ou la retirer. Vos preuves sont
conservées par l'école avec la décision.

\subsection{Le supplément émis}

La DEVE émet votre supplément en fin de cursus. Le PDF porte un \textbf{numéro} et un
\textbf{code de vérification}. Toute personne à qui vous le présentez -- un employeur, une
autre université -- peut vérifier, sans compte, sur la page \og Vérifier un supplément\fg{} de
l'école, qu'il a bien été émis et qu'il n'a pas été modifié. Si la DEVE en émet une nouvelle
version, l'ancienne est signalée comme remplacée.

\capture[0.8\linewidth]{gu-public-verification}{Vérification publique d'un supplément : numéro
et code reconnus.}

\newpage
% ===========================================================================
\part{Enseignant}\label{part:enseignant}
% ===========================================================================

\section{Vos rôles et votre page de suivi}

Un enseignant intervient à trois titres, qui se cumulent : \textbf{responsable} des EP du
catalogue qui lui sont confiés, \textbf{référent} des étudiants qu'il suit en stage,
\textbf{auteur} des EP qu'il propose. Vous ne voyez que ce qui relève de vous.

\subsection{Le suivi, toutes promotions confondues}

L'activité \og Suivi de l'enseignement personnalisé\fg{} de votre cours d'équipe est votre
point d'entrée : elle rassemble ce qui vous concerne dans toutes les promotions qui y sont
reliées. Le bouton \bouton{À traiter} affiche entre parenthèses le nombre de choses qui vous
attendent. Trois boutons mènent à ses pages :
\begin{description}
  \item[À traiter] tout ce qui attend une action de votre part : les EP dont vous êtes
    responsable avec des participants en attente (et l'ancienneté de la plus vieille demande),
    vos propositions et révisions de syllabus à reprendre, et les éléments du supplément à
    examiner comme référent, que vous tranchez directement sur la page.
  \item[Validation] toutes les demandes de votre périmètre, filtrables par promotion, type,
    statut, année et nom, avec un lien vers chacune.
  \item[Tableau de pilotage] l'avancement de vos étudiants référents face à leurs minimums.
\end{description}
La DEVE y trouve en plus un lien \chemin{Administration}, qui choisit les promotions dont les
demandes remontent ici (section~\ref{sec:synthese-admin}).

\capture{gu-ens-synthese-a-traiter}{Onglet \og À traiter\fg{} : participants à noter, syllabus,
éléments du supplément.}

\capture{gu-ens-synthese-validation}{Onglet \og Validation\fg{} : toutes les demandes de votre
périmètre.}

Chaque ligne renvoie vers la page de l'activité de la promotion concernée, décrite ci-dessous.

\section{Responsable d'un EP du catalogue}\label{sec:responsable}

\subsection{La liste des participants}

Dans l'activité de la promotion, \chemin{Validation} commence par \og Mes EP\fg{}. Le bouton
\bouton{Participants} d'un EP ouvre sa liste : en tête, son état, sa période d'inscription,
ses places, ses responsables et son barème ; puis un inscrit par ligne.

\capture{gu-ens-participants}{Participants d'un EP noté : une note par inscrit, enregistrée
pour toute la liste.}

\subsection{Noter ou valider toute la liste}

\begin{etapes}
  \item Pour un EP \textbf{noté}, saisissez la note de chaque inscrit (virgule ou point
    décimal) : la validation suit le seuil fixé ; une note sous le seuil vaut refus. Pour un EP
    \textbf{sans note}, choisissez \og Validé\fg{} ou \og Non validé\fg{}.
  \item Ajoutez au besoin un commentaire, transmis à l'étudiant.
  \item Cliquez sur \bouton{Enregistrer les décisions}.
\end{etapes}

Une ligne laissée vide n'est pas touchée. Le message de confirmation dit combien de décisions
ont été enregistrées, et combien de lignes avaient changé entre-temps (tranchées par quelqu'un
d'autre) : celles-là sont laissées telles quelles.

\begin{astuce}
Si le circuit du type a une étape après la vôtre (\og responsable puis DEVE\fg{}), votre
décision fait passer la demande à la DEVE. Sinon elle est définitive et l'étudiant est prévenu.
Une note portée à une étape précédente n'est montrée qu'en filigrane : elle ne décide de
l'étape courante que si vous la saisissez.
\end{astuce}

\subsection{Importer les notes}

\bouton{Importer les notes (CSV)} remplit toute la liste depuis un tableur.
\begin{etapes}
  \item Téléchargez la liste de l'EP (\bouton{Exporter la liste (CSV)}) : elle sert de modèle.
  \item Remplissez une ligne par étudiant, avec une ligne d'en-tête. Colonnes reconnues :
    \og Email\fg{} ou \og Numéro d'identification\fg{} ; \og Note\fg{} (EP noté) ou
    \og Décision\fg{} (\og validé\fg{}, \og oui\fg{}, \og 1\fg{} ou \og non validé\fg{},
    \og non\fg{}, \og 0\fg{}, \og absent\fg{}) ; \og Commentaire\fg{}, facultatif. Séparateur
    point-virgule, virgule ou tabulation.
  \item Déposez le fichier : la \textbf{prévisualisation} montre, ligne par ligne, ce qui sera
    fait (validé, refusé, corrigé, ignoré) et les erreurs (étudiant non inscrit, ligne en
    double, note illisible).
  \item \bouton{Confirmer l'import} applique les décisions, avec les mêmes contrôles qu'une
    saisie à la main.
\end{etapes}

\capture{gu-ens-import-notes}{Import des notes d'un EP depuis un fichier CSV.}

\subsection{Inscrire, retirer, clore}

\begin{itemize}
  \item \bouton{Inscrire des étudiants} : choisissez-les dans la liste (recherche par nom) ; ils
    sont inscrits en attente, même après la fermeture des inscriptions. Une inscription déjà en
    cours n'est pas doublée ; les places restent comptées.
  \item \bouton{Retirer} : retire un inscrit en attente, avec un motif ; la place est rendue.
  \item \bouton{Feuille d'émargement (PDF)} : la liste des inscrits à signer.
  \item \bouton{Clore l'EP} : plus aucune inscription, l'EP reste au catalogue et ses notes
    restent saisissables ; vous choisissez le sort des inscriptions encore en attente.
\end{itemize}

\section{Une demande à la fois}\label{sec:demande}

Le bouton \bouton{Valider} (ou \bouton{Détail}) d'une demande ouvre sa page : ce que
l'étudiant a demandé, ses justificatifs, l'étape en cours et l'historique.

\capture{gu-ens-demande}{Une demande en attente de votre décision : ECTS retenus, commentaire,
décision.}

\begin{itemize}
  \item \bouton{Valider} (\bouton{Valider définitivement} si vous êtes la dernière étape) :
    ajustez au besoin les \textbf{ECTS retenus} (au plus ceux demandés) ;
  \item pour un EP noté, saisissez la note puis \bouton{Enregistrer la note} : la validation
    suit le seuil ;
  \item \bouton{Refuser} : le commentaire est alors indispensable à l'étudiant.
\end{itemize}

Si quelqu'un a tranché la même demande entre-temps, rien n'est enregistré et un message vous
prévient que la demande a changé.

\section{Enseignant référent}\label{sec:referent}

\subsection{Les déclarations de vos étudiants}

Dans \chemin{Validation}, la rubrique \og En attente de votre décision\fg{} liste les
déclarations de vos étudiants arrivées à votre étape, de la plus ancienne à la plus récente.
En dessous, tout votre périmètre, filtrable.

\capture{gu-ens-validation}{Validation : ce qui attend votre décision, puis tout votre
périmètre.}

Pour chaque déclaration : ouvrez-la, lisez la description, ouvrez les justificatifs, puis
validez (avec les ECTS retenus) ou refusez en expliquant pourquoi (section~\ref{sec:demande}).

\subsection{Les éléments du supplément}

Les éléments que vos étudiants demandent pour leur supplément arrivent dans l'onglet
\og À traiter\fg{} du suivi et dans \chemin{Éléments de supplément à examiner} de l'activité :
catégorie, période, intitulés, description, preuves. \bouton{Valider} transmet l'élément à
l'étape suivante (en général la DEVE, qui décide en dernier), avec votre commentaire ;
\bouton{Refuser} l'arrête.

\capture{gu-ens-supplement-items}{Éléments du supplément à examiner, avec leurs preuves.}

\subsection{La relance du lundi}

Chaque lundi matin, un message rappelle à chaque validateur les demandes qui l'attendent depuis
plus d'une semaine.

\section{Proposer un EP, tenir son syllabus}\label{sec:auteur}

\subsection{Proposer un nouvel EP}

\begin{etapes}
  \item \chemin{Propositions et syllabus > Proposer un nouvel EP} : l'EP est créé à l'état de
    brouillon, et vous en devenez responsable.
  \item Renseignez ses caractéristiques : intitulé, type, ECTS, places, années concernées.
  \item Rédigez son \textbf{syllabus}, au modèle de la fiche ECTS : objectifs, acquis
    d'apprentissage, prérequis, contenu, méthodes, conditions d'examen et de seconde session,
    bibliographie, langue, travail personnel.
  \item Décrivez son \textbf{programme} : les séances, avec leur nature (cours, TD, TP,
    clinique\dots), leur durée et leur intervenant (un enseignant du site, ou un intervenant
    extérieur par son nom). Les séances ne sont pas datées : le syllabus décrit l'EP, pas le
    calendrier d'une session. Le volume horaire est calculé.
  \item \bouton{Enregistrer le brouillon} autant de fois que nécessaire, puis
    \bouton{Soumettre à la DEVE}.
\end{etapes}

\capture{gu-ens-proposer}{Proposition d'un nouvel EP : caractéristiques et syllabus.}

\capture{gu-ens-seances}{Programme des séances : intitulé, nature, durée, intervenant (compte du
site ou nom d'un intervenant extérieur).}

La DEVE publie la proposition (l'EP entre au catalogue), vous la renvoie avec un commentaire
(corrigez, puis soumettez de nouveau) ou la refuse. Vous êtes prévenu de sa décision, et la
proposition renvoyée apparaît dans l'onglet \og À traiter\fg{} du suivi.

\capture{gu-ens-propositions}{Propositions et syllabus : l'état de chaque version.}

\subsection{Mettre à jour un EP publié}

\bouton{Proposer une mise à jour} ouvre une \textbf{nouvelle version} du syllabus, préremplie :
la version en vigueur reste au catalogue tant que la nouvelle n'est pas publiée. Selon le
choix de la DEVE, la mise à jour est publiée aussitôt ou passe par sa revue. Les versions
précédentes restent consultables.

\capture{gu-ens-syllabus-edit}{Mise à jour du syllabus d'un EP publié : rubriques, puis
programme des séances.}

\subsection{Qui valide un syllabus}\label{sec:qui-valide-syllabus}

C'est la \textbf{DEVE} (capacité \texttt{mod/ep:manage}) qui valide les syllabus, depuis
\chemin{Propositions et syllabus} : elle publie la version soumise, la renvoie à son auteur avec
un commentaire, ou la refuse (section~\ref{sec:propositions}). Ni le référent ni un autre
enseignant n'interviennent.
\begin{center}
\begin{tabularx}{\linewidth}{@{}X l@{}}
\toprule
\textbf{Cas} & \textbf{Qui valide} \\
\midrule
Proposition d'un nouvel EP & Toujours la DEVE \\
Mise à jour du syllabus d'un EP publié, par son responsable & La DEVE ; ou personne, publiée
aussitôt, si la DEVE a coché la \og publication directe des mises à jour\fg{} \\
Syllabus ou séances importés par la DEVE (CSV, paramétrage JSON) & Publiés directement \\
\bottomrule
\end{tabularx}
\end{center}
Même publiée directement, une mise à jour reste dans l'historique des versions de l'EP.

\newpage
% ===========================================================================
\part{DEVE}\label{part:deve}
% ===========================================================================

\section{Le panneau d'administration}

Tout le paramétrage se fait dans l'activité de la promotion, depuis \chemin{Administration},
avec le rôle de la DEVE. Les pages y sont rangées par usage : règles d'attribution, catalogue,
éditions PDF et propositions, supplément au diplôme, imports et exports, suivi.

\capture{gu-deve-administration}{Panneau d'administration de l'activité.}

\begin{attention}
Par défaut, Moodle donne aussi les capacités de la DEVE au rôle \og enseignant éditeur\fg{}.
Si les enseignants ont ce rôle dans le cours de la promotion, retirez-leur
\texttt{mod/ep:validatedeve}, \texttt{mod/ep:manage} et \texttt{mod/ep:viewall} dans les
permissions de l'activité, et donnez à la DEVE le rôle \og gestionnaire\fg{}.
\end{attention}

\section{Mettre en place une promotion}\label{sec:miseenplace}

\begin{etapes}
  \item Dans le cours de la promotion, ajoutez une activité \og Enseignement personnalisé\fg{}
    et indiquez l'\textbf{année d'étude courante}.
  \item Si une autre promotion est déjà paramétrée, copiez ses règles :
    \chemin{Administration > Transférer le paramétrage > Copier depuis une autre activité}
    (section~\ref{sec:transfert}). Vérifiez, puis passez à l'étape 6.
  \item Réglez les \textbf{types} et leurs circuits (section~\ref{sec:types}).
  \item Fixez les \textbf{minimums} (section~\ref{sec:minimums}).
  \item Réglez le \textbf{supplément au diplôme} (section~\ref{sec:supplement-reglages}) et les
    \textbf{courriels} (section~\ref{sec:courriels}).
  \item Constituez le \textbf{catalogue} (section~\ref{sec:catalogue}) et publiez-le.
  \item Dans le cours d'équipe, reliez l'activité de \og Suivi\fg{} à la nouvelle promotion
    (section~\ref{sec:synthese-admin}).
\end{etapes}

\section{Types d'EP et circuits de validation}\label{sec:types}

\chemin{Administration > Types d'EP et plafonds} liste les types de la promotion, avec leur
mode d'acquisition, leur circuit, leur évaluation, leurs justificatifs et leurs plafonds.
Vous pouvez en ajouter, les modifier, les réordonner, les désactiver ; un type qui porte déjà
des EP ne se supprime pas.

\capture{gu-deve-types}{Types d'EP de la promotion.}

\capture{gu-deve-type-edit}{Réglage d'un type : acquisition, circuit, évaluation, justificatifs,
plafonds.}

Pour chaque type :
\begin{description}
  \item[Mode d'acquisition] inscription à un EP du catalogue ; déclaration par l'étudiant ;
    automatique, d'après les stages complémentaires validés (avec un nombre d'ECTS par jour de
    stage) ; ou saisie par la DEVE seule, à la main ou par import.
  \item[Circuit de validation] une à trois étapes, dans l'ordre choisi parmi : responsable de
    l'EP, enseignant référent, DEVE ; ou \og aucune\fg{} : validé d'office. La DEVE peut
    toujours trancher à la place de l'étape en cours.
  \item[Évaluation] sans note (validé / non validé), ou notée avec une note maximale et un seuil
    de validation. Un EP du catalogue peut déroger à la règle de son type.
  \item[Justificatifs] sans objet, facultatifs ou obligatoires, et les formats admis.
  \item[Plafonds] maximum d'ECTS retenus sur le cursus et par année.
  \item[Consigne] texte affiché à l'étudiant quand il déclare un EP de ce type.
  \item[Supplément] le type figure ou non à la rubrique 4.3 du supplément.
\end{description}

\begin{astuce}
Un changement de circuit vaut pour les nouvelles demandes et pour les étapes restantes des
demandes en cours ; les décisions déjà prises restent dans l'historique. Le bouton
\bouton{Resynchroniser les EP de type stage} recalcule aussitôt les ECTS attribués d'après les
stages ; il n'est utile que pour vérifier le résultat, car le calcul est refait de lui-même dès
qu'un stage ou le barème change : la DEVE peut valider un stage dans « Gestion des stages » puis,
aussitôt, les années ou le supplément, le stage y compte.
\end{astuce}

\section{Minimums d'ECTS}\label{sec:minimums}

\chemin{Administration > Minimums d'ECTS} fixe le minimum d'ECTS à valider pour chaque année
d'étude (une année laissée à 0 n'impose rien) et, en plus, un minimum sur l'ensemble du
cursus. Les tableaux de bord, le pilotage et la validation des années s'y réfèrent.

\capture{gu-deve-minimums}{Minimums d'ECTS par année et sur le cursus.}

\section{Le catalogue}\label{sec:catalogue}

\subsection{Gérer les EP}

\chemin{Administration > Catalogue des EP internes} liste les EP : type, ECTS, années, places
restantes, inscrits, responsables, état ; avec, pour chacun, \bouton{Modifier},
\bouton{Responsables}, \bouton{Participants} et les changements d'état.

\capture{gu-deve-catalogue}{Catalogue des EP vu par la DEVE.}

La fiche d'un EP (\bouton{Ajouter un EP au catalogue} ou \bouton{Modifier}) réunit :
\begin{itemize}
  \item description, type, ECTS, années minimale et maximale, nombre de places, code, ordre
    d'affichage ;
  \item \textbf{état} et périodes d'\textbf{inscription} et de \textbf{déroulement} ;
  \item \textbf{évaluation} (celle du type, ou une dérogation) ;
  \item \textbf{prérequis} : EP à avoir validés, ECTS minimum, bloquants ou non.
\end{itemize}

\capture{gu-deve-ep-edit}{Fiche d'un EP du catalogue.}

\bouton{Responsables} désigne les enseignants qui valideront les inscriptions.

\capture[0.8\linewidth]{gu-deve-responsables}{Responsables d'un EP.}

Les \textbf{états} d'un EP :
\begin{itemize}
  \item \statut{Brouillon}, \statut{Proposé} : en préparation, invisible des étudiants ;
  \item \statut{Publié} : au catalogue ; inscriptions ouvertes pendant leur période ;
  \item \statut{Clos} : plus d'inscription ; les inscrits restent et sont notés ;
  \item \statut{Archivé} : hors du catalogue ; l'historique reste.
\end{itemize}

Le \bouton{Livret des EP (PDF)} rassemble les fiches de tous les EP publiés ou clos, précédées
d'un sommaire.

\subsection{Importer le catalogue}

\chemin{Administration > Importer le catalogue (CSV)} crée ou met à jour des EP, avec leurs
responsables, leurs prérequis et leur syllabus ; \chemin{Importer les séances des syllabus
(CSV)} remplace leur programme. Chaque import montre d'abord une prévisualisation, puis
s'applique sur confirmation ; un modèle se télécharge depuis la page. Les colonnes sont
décrites dans le document \emph{Formats des imports et des exports}.

\capture{gu-deve-import-catalogue}{Import du catalogue depuis un fichier CSV.}

\subsection{Examiner les propositions des enseignants}\label{sec:propositions}

\chemin{Propositions et syllabus} liste les versions soumises par les enseignants. Pour
chacune, après lecture de la fiche (et, pour une mise à jour, la comparaison avec la version
en vigueur) : \bouton{Publier}, \bouton{Renvoyer à l'auteur} avec un commentaire, ou
\bouton{Refuser}. L'auteur est prévenu.

\capture{gu-deve-propositions}{Propositions et mises à jour à examiner.}

\capture{gu-deve-comparer}{Mise à jour soumise : version en vigueur, version proposée et décision.}

\section{Éditions PDF et courriels}\label{sec:courriels}

\chemin{Éditions PDF et propositions} règle l'en-tête, le pied de page et le logo des fiches,
du livret, des feuilles d'émargement et du supplément, et permet aux responsables de publier
directement les mises à jour de leurs syllabus (une proposition de nouvel EP passe toujours par
la DEVE).

\chemin{Courriels} règle le sujet et le texte de chaque message (section~\ref{sec:messages}),
et l'active ou non. Les variables entre accolades (\texttt{\{student\}}, \texttt{\{ep\}},
\texttt{\{decision\}}, \texttt{\{ects\}}, \texttt{\{grade\}}, \texttt{\{comment\}},
\texttt{\{link\}}\dots{}) sont remplacées à l'envoi ; celles de chaque message sont rappelées
sous son texte. Un message laissé vide reprend le texte par défaut.

\capture{gu-deve-courriels}{Texte des courriels et leurs variables.}

\section{Suivre la promotion}\label{sec:suivi}

\subsection{Validation}

\chemin{Validation} montre d'abord les demandes en attente (40 par page), puis toute la
promotion, filtrable par type, statut, année d'étude et nom. La DEVE peut :
\begin{itemize}
  \item trancher à n'importe quelle étape, à la place de l'étape en cours ;
  \item \textbf{corriger} une décision finale (ECTS retenus, note), avec un commentaire gardé
    dans l'historique ;
  \item \textbf{retirer} une demande, même tranchée (\bouton{Retirer cet EP}), avec un
    commentaire.
\end{itemize}

\capture{gu-deve-validation}{Validation, vue de la DEVE : toute la promotion.}

\subsection{Tableau de pilotage}

\chemin{Tableau de pilotage} présente l'avancement de chaque étudiant face à ses minimums,
filtrable, avec un lien vers sa situation détaillée (le même tableau de bord que celui de
l'étudiant).

\capture{gu-deve-pilotage}{Tableau de pilotage de la promotion.}

\subsection{Export CSV}

\chemin{Export CSV} produit le \bouton{Bilan par étudiant} (une ligne par étudiant : ECTS par
année et par type, minimums atteints) ou le \bouton{Détail des EP} (une ligne par EP de chaque
étudiant, avec son statut et ses décisions), à ouvrir dans un tableur.

\section{Fin d'année}\label{sec:finannee}

\subsection{Valider les années}

\chemin{Administration > Validation des années} présente, pour l'année d'étude choisie, le
bilan de chaque étudiant face au minimum : ECTS retenus, minimum atteint ou non, décision déjà
prise.

\capture{gu-deve-annees}{Validation des années d'étude.}

\begin{etapes}
  \item \bouton{Valider les années atteintes} : enregistre en une fois la validation de tous
    les étudiants qui ont atteint le minimum. Le message dit combien d'années ont été
    validées, et combien sont restées sans décision.
  \item Pour les autres, décidez étudiant par étudiant, après le jury : validée (dérogation,
    avec un commentaire) ou non validée ; ou, si le jury le décide, \bouton{Déclarer non validées
    les années non atteintes} en une fois.
\end{etapes}

Les décisions d'année figurent au supplément au diplôme.

\subsection{Rentrée}

\chemin{Administration > Rentrée} fait passer l'année courante de l'activité à la suivante :
les objectifs de la nouvelle année deviennent exigibles ; les crédits, décisions et suppléments
restent. La page rappelle d'abord combien d'étudiants n'ont pas d'année validée.

La même page \textbf{transfère} le dossier d'un étudiant (EP, éléments du supplément et
preuves, décisions d'année) vers l'activité d'une autre promotion, par exemple pour un
redoublement. L'étudiant doit d'abord être inscrit dans le cours de destination.

\capture{gu-deve-rentree}{Rentrée et transfert d'un étudiant.}

\section{Le supplément au diplôme}\label{sec:supplement}

\subsection{Paramétrage}\label{sec:supplement-reglages}

\chemin{Administration > Supplément au diplôme : paramétrage} :
\begin{itemize}
  \item le \textbf{circuit} des éléments complémentaires (ajout direct, référent, DEVE, ou
    référent puis DEVE) et l'exigence de \textbf{preuves} ;
  \item le \textbf{préfixe} des numéros, le champ du profil qui porte la date de naissance, le
    \textbf{cachet} ;
  \item les \textbf{textes} des rubriques fixes du modèle Europass, en français et en anglais :
    qualification, niveau, conditions d'accès, programme, fonction du diplôme, système
    national, \textbf{signataire} et sa qualité ;
  \item les \textbf{catégories} d'éléments complémentaires.
\end{itemize}

\capture{gu-deve-supplement-reglages}{Paramétrage du supplément au diplôme.}

\subsection{Éléments à examiner}

\chemin{Éléments de supplément à examiner} présente les demandes arrivées à l'étape de la
DEVE, avec leurs preuves et l'avis du référent, à valider ou à refuser directement dans la
liste ; en dessous, tous les éléments de la promotion, filtrables par statut et par nom.

\capture{gu-deve-supplement-elements}{Éléments du supplément : file d'attente de la DEVE.}

\subsection{Émission}

\chemin{Administration > Émission des suppléments} liste les étudiants, ce qui manque encore
à leur supplément (date de naissance, numéro d'identification, signataire) et leur supplément
en vigueur.

\capture{gu-deve-emission}{Émission des suppléments au diplôme.}

\begin{itemize}
  \item Pour un étudiant : \bouton{Aperçu}, puis \bouton{Émettre le supplément}, en français
    ou en anglais ; \bouton{Émettre malgré les informations manquantes} si besoin.
  \item Pour toute la promotion : \bouton{Émettre pour toute la promotion}. L'émission tourne
    en tâche de fond ; un message vous prévient à la fin, avec le nombre de suppléments émis et
    non émis. Sur le banc, les 1\,000 suppléments d'une promotion sont émis en quelques
    minutes.
  \item Une \textbf{archive ZIP} rassemble les PDF émis.
\end{itemize}

Chaque supplément émis porte un numéro unique et un code de vérification ; il est conservé ;
une nouvelle émission remplace la précédente, qui reste vérifiable comme \og remplacée\fg{}.
\chemin{Vérifier un supplément}, page accessible sans compte, confirme qu'un document présenté
a bien été émis et n'a pas été modifié.

\section{Imports, exports, transfert}\label{sec:transfert}

\begin{description}
  \item[Transférer le paramétrage] exporter les règles (types, minimums, supplément, courriels,
    catalogue et syllabus) dans un fichier JSON, importer un tel fichier partie par partie, ou
    copier directement les règles d'une autre activité que vous gérez ; le journal des imports
    garde la trace de chacun.
  \item[Importer des crédits (CSV)] porter au dossier des étudiants des EP suivis ailleurs, ou
    l'historique d'un outil précédent, avec prévisualisation.
  \item[Export complet et restauration] télécharger toute l'activité dans une archive ZIP
    (classeur Excel lisible, paramétrage, données des étudiants, fichiers), ou restaurer une
    telle archive dans une autre activité : ce qui existe déjà est écarté, rien n'est supprimé.
\end{description}

\capture{gu-deve-transfert}{Transfert du paramétrage et journal des imports.}

\capture{gu-deve-archive}{Export complet et restauration.}

La sauvegarde Moodle du cours reste l'outil de référence pour conserver une promotion ; la
réinitialisation du cours efface, au choix, les dossiers des étudiants, les décisions d'année,
les suppléments émis et le journal des imports, en gardant le paramétrage et le catalogue.

\section{Le suivi des enseignants}\label{sec:synthese-admin}

Dans le cours d'équipe, ajoutez une activité \og Suivi de l'enseignement personnalisé\fg{}, puis,
dans son onglet \chemin{Administration}, cochez les activités \og Enseignement personnalisé\fg{}
des promotions à faire remonter. Une promotion qui n'est plus suivie se décoche sans toucher à
son activité. Les droits restent ceux de chaque activité de promotion.

\capture{06-admin-lier-synthese}{Choix des promotions qui remontent dans le suivi.}

\newpage
% ===========================================================================
\part{Administrateur du site}\label{part:admin}
% ===========================================================================

\section{Installation et exploitation}

L'installation, la mise à niveau, les tâches planifiées, les capacités et le réglage du
serveur sont décrits dans \texttt{mod/ep/INSTALL.md}. L'essentiel :
\begin{itemize}
  \item Moodle 4.3 ou plus récent, PHP 8.0 ou plus récent ; \texttt{mod\_stage} installé ;
  \item aucun réglage d'administration du site : tout se fait dans l'activité ;
  \item trois tâches : recalcul nocturne des EP de type stage, relance du lundi, émission
    groupée des suppléments (tâche ad hoc : le \emph{cron} doit tourner) ;
  \item pour une école de plusieurs milliers d'étudiants : connexions persistantes à la base,
    APCu pour le cache applicatif, JIT et OPcache de PHP (gains mesurés dans
    \texttt{INSTALL.md}).
\end{itemize}

\begin{center}
\begin{tabularx}{\linewidth}{@{}l X@{}}
\toprule
\textbf{Capacité} & \textbf{Pour qui (par défaut)} \\
\midrule
\texttt{mod/ep:view} & Tous les inscrits \\
\texttt{mod/ep:submit} & Étudiants : s'inscrire, déclarer, demander des éléments du
supplément \\
\texttt{mod/ep:evaluateteacher} & Enseignants : trancher ce dont ils ont la charge \\
\texttt{mod/ep:propose} & Enseignants : proposer un EP, mettre à jour le syllabus du leur \\
\texttt{mod/ep:validatedeve} & DEVE : trancher et corriger pour tous, valider les années \\
\texttt{mod/ep:manage} & DEVE : paramétrage, catalogue, supplément, imports \\
\texttt{mod/ep:viewall} & DEVE : voir tous les étudiants, exporter \\
\bottomrule
\end{tabularx}
\end{center}

\newpage
% ===========================================================================
\part{Questions fréquentes}\label{part:faq}
% ===========================================================================

\section{Étudiants}

\begin{description}
  \item[Le bouton \og S'inscrire\fg{} n'apparaît pas.] La raison est écrite à sa place :
    inscriptions fermées, EP complet, année non concernée, déjà inscrit, prérequis manquants
    (section~\ref{sec:inscription}).
  \item[Je me suis inscrit au mauvais EP.] Tant que l'inscription est en attente,
    \bouton{Retirer ma demande} sur le tableau de bord.
  \item[Mes ECTS validés ne comptent pas tous.] Le plafond du type est atteint : le bilan par
    type les montre en \og non retenus\fg{}.
  \item[La page de déclaration dit que je n'ai pas de référent.] Il n'est pas encore attribué
    dans \og Gestion des stages\fg{} : prévenez la DEVE.
  \item[Mon stage n'apparaît pas.] Seuls les stages \emph{complémentaires} validés par la DEVE
    donnent des ECTS d'EP, et seulement si la DEVE a fixé un nombre d'ECTS par jour.
  \item[L'aperçu du supplément signale des informations manquantes.] Date de naissance ou
    numéro d'identification : à faire corriger dans votre profil par l'école ; signataire : la
    DEVE le renseignera avant l'émission.
\end{description}

\section{Enseignants}

\begin{description}
  \item[Je ne vois pas un étudiant que je suis en stage.] Vérifiez dans \og Gestion des
    stages\fg{} que vous êtes bien son enseignant référent : c'est cette attribution qui compte.
  \item[Ma liste de participants dit que des lignes ont changé.] Quelqu'un (la DEVE, un autre
    responsable) a tranché ces demandes pendant votre saisie : elles n'ont pas été écrasées.
  \item[Mon import de notes signale des lignes.] Étudiant non inscrit à l'EP, ligne en double,
    note hors barème ou décision non reconnue : la prévisualisation les liste, corrigez le
    fichier et déposez-le de nouveau.
  \item[Ma mise à jour de syllabus n'apparaît pas au catalogue.] Elle attend la DEVE : l'onglet
    \og À traiter\fg{} du suivi la montre \og En attente de la DEVE\fg{}.
\end{description}

\section{DEVE}

\begin{description}
  \item[Un enseignant voit tout le monde.] Il a le rôle \og enseignant éditeur\fg{} avec les
    capacités de la DEVE : retirez-les (section~\ref{sec:miseenplace}).
  \item[Les EP de type stage restent à 0.] Le type stage n'a pas d'ECTS par jour : réglez-le,
    puis \bouton{Resynchroniser les EP de type stage}.
  \item[L'émission groupée ne se termine pas.] Elle tourne en tâche de fond : le \emph{cron}
    du site doit tourner. Le message de fin indique les suppléments non émis faute
    d'informations.
  \item[Je veux reprendre les règles de l'an dernier.] \chemin{Transférer le paramétrage >
    Copier depuis une autre activité}.
  \item[Un étudiant redouble.] \chemin{Rentrée > Transférer un étudiant} vers l'activité de sa
    nouvelle promotion, après l'avoir inscrit dans le cours correspondant.
\end{description}

\appendix
\section{Aide-mémoire des pages}

\begin{center}
\small
\begin{tabularx}{\linewidth}{@{}l l X@{}}
\toprule
\textbf{Page} & \textbf{Qui} & \textbf{Pour} \\
\midrule
Mes enseignements personnalisés & Étudiant & Bilan, détail, retrait \\
Catalogue des EP & Étudiant, tous & Rechercher, lire les syllabus, s'inscrire \\
Déclarer un EP & Étudiant & Déclarer avec justificatifs \\
Mon supplément au diplôme & Étudiant & Aperçu, éléments, supplément émis \\
Validation & Enseignant, DEVE & Trancher, corriger, filtrer \\
Participants & Responsable, DEVE & Noter en lot, importer, inscrire, retirer, clore \\
Propositions et syllabus & Enseignant, DEVE & Proposer, mettre à jour, examiner \\
Éléments de supplément à examiner & Référent, DEVE & Trancher les éléments \\
Tableau de pilotage & Enseignant, DEVE & Avancement face aux minimums \\
Suivi : À traiter, Validation & Enseignant & Tout ce qui attend, toutes promotions \\
Administration & DEVE & Tout le paramétrage (partie~\ref{part:deve}) \\
Vérifier un supplément & Tous, sans compte & Authenticité d'un supplément \\
\bottomrule
\end{tabularx}
\end{center}

\section{Captures}

Les captures se refont sur le banc d'essai (\texttt{docs/audit/banc}) avec
\texttt{docs/guides/captures.js} : \texttt{node captures.js site-<promotion>.json gu-...},
avec les variables \texttt{PERSONNEL}, \texttt{ETU}, \texttt{RESP}, \texttt{VERIF} et
\texttt{SYNTHESE} décrites en tête du script. Les captures des propositions de syllabus et du
suivi viennent des procédures de bout en bout du banc.

\end{document}
